% =====================================================================
%  c00.tex —— 标题页 + 译者说明 + 目录（Lead 独占）
% =====================================================================
\begin{titlepage}
\centering
{\small\ttfamily IFT-UAM/CSIC-21-151, ZMP-HH/21-28\par}
\vspace*{1.6cm}

{\LARGE\bfseries 4d $\boldsymbol{\mathcal{N}=2}$ 超共形场论漫游指南\par}

\vspace{1.0cm}
{\large\itshape The Hitchhiker's Guide to\\[0.2em]
4d $\mathcal{N}=2$ Superconformal Field Theories\par}

\vspace{1.4cm}
{\large Mohammad Akhond,\ Guillermo Arias-Tamargo,\ Alessandro Mininno,\\[0.3em]
Hao-Yu Sun,\ Zhengdi Sun,\ Yifan Wang,\ Fengjun Xu\par}

\vspace{0.6cm}
{\normalsize arXiv:2112.14764v3 [hep-th]（2022 年 8 月 21 日）\par}

\vspace{1.2cm}
\begin{minipage}{0.86\linewidth}
\small
\noindent\textbf{摘要}\quad 四维时空中带 $\cN=2$ 超对称的超共形场论，是研究量子场论中强耦合现象的一块绝佳试验场。
其刚性的结构保证了对非微扰效应的宝贵解析控制，而理论本身又足够灵活，足以容纳一个庞大的量子系统图景。
本文旨在提供一部导引手册，介绍四维 $\cN=2$ 超共形场论的基本特征，以及在弦/M/F 理论中构造它们的基本工具。
内容基于 2020 年 7 月「量子场论与几何」暑期学校（Quantum Field Theories and Geometry School）上的系列讲座。
\end{minipage}

\vfill
{\normalsize 中 文 译 注 本\par}
\vspace{0.4cm}
{\small 逐段中译 + 学术批注（\texttt{\% [Tutor Note: ...]}）；公式编号与原文一一对应\par}
\vspace{1.0cm}
\end{titlepage}

\section*{译者说明}
\addcontentsline{toc}{section}{译者说明}

\noindent\textbf{底本。}本译注本的底本是 M.~Akhond, G.~Arias-Tamargo, A.~Mininno,
H.-Y.~Sun, Z.~Sun, Y.~Wang, F.~Xu 合著的讲义
\emph{The Hitchhiker's Guide to 4d $\mathcal{N}=2$ Superconformal Field Theories}
（arXiv:2112.14764v3 [hep-th]，2022 年 8 月 21 日，共 122 页），内容基于 2020 年 7 月
Quantum Field Theories and Geometry School 的系列讲座。底本 PDF 页号为 $1$--$122$，
印刷页码自 PDF 第 3 页起为 $2$，即 \textbf{PDF 页号 $=$ 印刷页码 $+1$}；
全书正文止于 PDF p.97（致谢），参考文献自 PDF p.97 起至 p.122。

\medskip
\noindent\textbf{体例。}
\begin{itemize}[leftmargin=1.6em,itemsep=2pt]
\item 正文逐段全译。原文的公式编号用 \verb|\tag{}| 原样复刻（如 (3.13)），
      共 164 条，与原文一一对应，不缺、不重、不跳、不自造；
\item 每段原文页起始处插注释 \verb|% [原文 p.N]|，便于与底本对页；
\item 原文的定义（Definition）、例（Example）、习题（Exercise）按原版的左边框／方框样式重排，
      编号照抄原文（定义 2.1--2.5；例 3.1、4.1--4.9、5.1--5.3；习题 2.1--2.4、3.1--3.2、4.1--4.3、5.1）；
\item 15 幅插图按原文矢量图 450\,dpi 截取（图内英文标注为原图固有，不再重绘），图注中译；
\item 8 张表格按 \texttt{booktabs} 重排，表题中译，编号照抄原文；
\item 解释性补充、易错点提示、原文笔误的判定，一律写在 \verb|% [Tutor Note: ...]| 批注里，
      \textbf{不进入排版}；全部批注另见随包交付的《译注汇总.md》；
\item 参考文献不翻译（学术惯例），保留英文原貌并声明为自动转录、未逐条校对。
\end{itemize}

\medskip
\noindent\textbf{提示。}本译本尽量做到「正文一字不改」：原文的笔误照原样排印，
只在批注中指出并给出正确形式。若读者发现译文与底本不符，请以
arXiv:2112.14764v3 为准。

\clearpage
\tableofcontents
\clearpage
